\clearpage\thispagestyle{empty}
\pdfbookmark[0]{Quellen zur Anzeige über Picot und Sigorgne}{picot1808-sources}
\begin{center}\textbf{Quellen zur Anzeige über Picot und Sigorgne}\end{center}
\noindent Diese KI-Übertragung ins heutige Deutsch folgt Gauß’ Anzeige auf S.~507 der \href{https://gdz.sub.uni-goettingen.de/id/PPN236007467}{\textit{Werke}, Bd.~VI} (Göttingen, 1874). Sie erschien ohne Namensnennung unter der Überschrift \emph{Genf} in den \textit{Göttingischen gelehrten Anzeigen}, Stück~164, 13. Oktober 1808, S.~\href{https://api.digitale-sammlungen.de/iiif/image/v2/bsb10538674_00623/full/1800,/0/default.jpg}{1639}--\href{https://api.digitale-sammlungen.de/iiif/image/v2/bsb10538674_00624/full/1800,/0/default.jpg}{1640}. Das verlinkte Exemplar stammt aus der Bayerischen Staatsbibliothek.
\par\smallskip\noindent Die französischen Titel in der Anzeige bedeuten \emph{Über die einundzwanzig jüngsten Kometen und die neuen Planeten}, von Professor P. Picot, und \emph{Beweis der unmittelbaren Erschaffung der Erde im festen Zustand und der Unmöglichkeit physikalischer Ursachen für die Ausbildung ihrer Gestalt}, von Abbé Sigorgne. Im ersten Beitrag berichtet Gauß über den Auszug aus einem Brief von Olbers. KI-Übertragung und KI-Satz.
\clearpage
